% Includable, self-contained research appendix. Requires amsmath, amssymb,
% amsthm and theorem/lemma/proposition/corollary environments.
% All new labels begin zcc:. Mathematical source and license are attributed below.
\section{A narrower congruence rule and its exact boundary}
\label{zcc:section}

This section strengthens the supplied report \emph{Counter Congruence
Suffices for iPRA}, specifically its \texttt{theory/counter\_theorem.tex},
in \texttt{ipra\_counter\_congruence\_2026-11-10.zip} at ProveIt commit
\texttt{76a2651601d17b34d46753f189ea026a21ce8d29}. The numerical basis,
paired logical calculus, root-replacement equality algebra, and formula
lifting are inherited. The refinement is that counter congruence needs
only a literal-zero seed and updates independent of the changing state.
A shorter representation reduces its numerical-congruence proof overhead.
A second, distinct rule permits an even shorter derivation. Finally a
rule countermodel shows why restricting the premise endpoints to
successors is not the same simplification.

These are mathematical inference results. No smaller checker implementing
these bases is claimed. The supplied source and documentation use the
MIT No Attribution License, copyright 2026 ProveIt Contributors; that
license is retained with the unchanged reference checker and producer.
Mathematical attribution is retained regardless of the license's permission.

\subsection{Complete calculus used by the refinement}

There is one numerical sort. Finite terms are generated by variables,
$0$, $S$, and $R(n,a;u,v.h)$. Only the body $h$ binds $u,v$: $u$ is the
old counter and $v$ is the previous state. Other variables are passive
parameters. Successor may instead be an abbreviation
$Sa:=R(0,a;u,v.0)$, with fresh unused $u,v$. For a conventional first-order
signature, introduce one function symbol per fixed, finite earlier-stage
body description; no variable ranges over functions or descriptions, and
there is no uniformly total evaluator. The computation schemas are
\[
 R(0,a;u,v.h)=Sa,\qquad
 R(Sn,a;u,v.h)=h[n/u,R(n,a;u,v.h)/v].
 \tag{R0, RS}\label{zcc:computation}
\]
The substitutions are simultaneous and capture-free. In particular,
the scalar-state iterator from the smallest measured checker is not being
silently used here: this $R$ body receives \emph{both} counter and state.

Object formulas and internal resources have the grammars
\[
 A::=(s=t)\mid\mathsf D_{x,y}A,
 \qquad \mathsf Q_yA(y).
\]
The variables $x,y$ of $\mathsf D$ are distinct. Contexts $\Gamma$ are
finite reusable lists of object assumptions and internal resources.
An object hypothesis has the one-node rule $\Gamma\vdash A$ when
$A\in\Gamma$. An internal lookup, written
$\Gamma\mathrel{\triangleright}\mathsf Q_yB$, also has one node when
that resource belongs to $\Gamma$; it is not an object-formula conclusion.
R0 and RS are one-node object axioms in every retained context.
The logical rules are exactly
\[
 \frac{\Gamma\mathrel{\triangleright}\mathsf Q_yB(y)
       \qquad\Gamma\vdash B(0)}
      {\Gamma\vdash B(t)}\quad(\mathsf{QE}),
\]
\[
 \frac{\Gamma\vdash\mathsf D_{x,y}B(x,y)\qquad
       \Gamma,\mathsf Q_yB(a,y)\vdash F}
      {\Gamma\vdash F}\quad(\mathsf{DE}),
\]
\[
 \frac{\Gamma,B(t,0)\vdash B(t,Sz)}
      {\Gamma\vdash\mathsf D_{x,y}B(x,y)}
      \quad(\mathsf{DI}_{0}),
 \qquad
 \frac{\Gamma,E(t,z)\vdash E(t,Sz)}
      {\Gamma\vdash\mathsf D_{x,y}E(x,y)}
      \quad(\mathsf{DI}_{1}).
\]
The progress matrix $E$ must be one literal equation, whereas $B$ may
be any object formula. In DE, $a$ is fresh for the retained context,
conclusion, and fixed parameters of $B$. In either introduction, $z$
is fresh for the context, witness $t$, and fixed matrix parameters.
Local assumptions are discharged. Weakening, exchange, contraction,
renaming, and capture-free numerical substitution are admissible
structural operations, not additional proof nodes.

The equality rule retained in every version below is root-left replacement:
\[
 \frac{\Gamma\vdash s=t\qquad\Gamma\vdash s=r}
      {\Gamma\vdash t=r}\quad(\mathsf{RL}).
 \label{zcc:root}
\]
Both displayed premises may be arbitrary proofs. Restricting its major
to a retained hypothesis is \emph{not} part of the result.
The proposed congruence rule is
\[
 \frac{\Gamma\vdash s=t}
 {\Gamma\vdash R(s,0;u,v.h(u,\vec p))=
                  R(t,0;u,v.h(u,\vec p))}
 \quad(\mathsf{ZCC}),
 \label{zcc:rule}
\]
where $v$ has no free occurrence in the displayed body, including below
nested numerical binders. Its seed is literally $0$, not merely a term
proved equal to zero. There are no primitive reflexivity, general
congruence, transport, coverage, separation, or additional induction rules.
Call the resulting calculus $\mathcal C_0$.

\subsection{Equality and cases before using ZCC}

\begin{lemma}[Root equality algebra]\label{zcc:algebra}
RL and computation derive reflexivity, symmetry and transitivity.
From proofs of $L=R$, $L=F$, and $R=G$, exactly two additional RL nodes
produce a proof of $F=G$.
\end{lemma}
\begin{proof}
Set $c_s=R(Ss,0;u,v.u)$. RS gives $c_s=s$. Using two copies as the
major and minor of RL gives $s=s$, in three nodes. An available proof
of $s=t$ and this reflexivity give $t=s$. For transitivity, reverse
$s=t$ to $t=s$ and apply RL with minor $t=r$.
For the last statement, RL applied to $L=R$ and $L=F$ gives $R=F$.
Apply RL again with that major and minor $R=G$ to obtain $F=G$.
\end{proof}

\begin{lemma}[Numerical cases]\label{zcc:cases}
For every object template $A(y)$, from $\Gamma\vdash A(0)$ and
$\Gamma\vdash A(Sz)$ one can derive $\Gamma\vdash A(t)$, where $z$
is fresh for the retained data and $t$. The construction uses neither
equality rule nor progress introduction and adds four nodes to the
sum of the two premise trees.
\end{lemma}
\begin{proof}
Choose $x$ unused in $A$. Weaken the successor proof by $A(0)$ and apply
case introduction with witness $0$, producing $\mathsf D_{x,y}A(y)$.
Eliminate it, obtaining the resource $\mathsf Q_yA(y)$ because the
witness coordinate is unused. Look up that resource and use QE with
the original base proof, weakened to this context. Close DE. The added
nodes are DI, DE, resource lookup, and QE.
\end{proof}

\subsection{Zero seed and no state dependency suffice}

Define the fixed numerical context
\[
 U(n):=R(n,0;u,v.S(Su)).
 \label{zcc:U}
\]
Its two direct computations are
\[
 U(0)=S0,\qquad U(Sz)=S(Sz).
\]
Numerical cases therefore prove
\[
                         U(n)=Sn
 \tag{U}\label{zcc:Urep}
\]
for arbitrary $n$, with arbitrary retained context. This proof has six
nodes and uses no congruence or induction rule.

\begin{theorem}[Numerical congruence from ZCC]\label{zcc:numerical}
For every fixed numerical template $F(x,\vec p)$, $\mathcal C_0$ admits
\[
 \frac{\Gamma\vdash s=t}
      {\Gamma\vdash F(s,\vec p)=F(t,\vec p)}.
\]
Its displayed construction adds exactly twenty nodes to the supplied
equality proof and introduces no progress inference.
\end{theorem}
\begin{proof}
Freshen all new binders away from the passive parameters and source terms.
Apply ZCC to $s=t$ through $U$. It is permitted because $U$ has seed zero
and its update $S(Su)$ ignores $v$. Thus $U(s)=U(t)$. The two representation
proofs \eqref{zcc:Urep} give $U(s)=Ss$ and $U(t)=St$.
Two root replacements, as in Lemma~\ref{zcc:algebra}, give
\[
                              Ss=St.
 \label{zcc:successor}
\]

For the arbitrary template $F$, set
\[
                  V_F(n):=R(Sn,0;u,v.F(u,\vec p)).
 \label{zcc:V}
\]
Its update ignores its state. A single RS instance proves
\[
                         V_F(n)=F(n,\vec p)
 \tag{V}\label{zcc:Vrep}
\]
for every $n$. No cases or equality substitution are hidden in this
computation. Apply ZCC to the derived $Ss=St$ with body $F(u,\vec p)$;
its conclusion is $V_F(s)=V_F(t)$. The two endpoint computations
\eqref{zcc:Vrep}, followed by two more root replacements, give
$F(s,\vec p)=F(t,\vec p)$.

For a source proof with $P$ nodes, the successor stage has
$(P+1)+6+6+2=P+15$ nodes: one ZCC, two six-node representations, and two
RL nodes. The final stage adds one ZCC, two one-node computations, and
two RL nodes. The total is $P+20$. Each supplied premise occurs once.
\end{proof}

The supplied report's displayed numerical-congruence construction has
$P+56$ nodes. Merely replacing its final symmetry/transitivity assembly
by the two-root-replacement lemma would improve that existing construction
to $P+44$. The present zero-seed representation gives $P+20$. Thus the
comparison separates the new representation from a shorter assembly of
old proofs. These are exact sizes of specified proof trees, not minimum
proof-size claims. Their term labels contain $F$ and may grow with its size.

\subsection{Formula transport and full consequence equivalence}

\begin{theorem}[Full transport]\label{zcc:formula}
For every native Eq/$\mathsf D$ template $A(w)$, $\mathcal C_0$ admits
\[
 \frac{\Gamma\vdash s=t\qquad\Gamma\vdash A(s)}
      {\Gamma\vdash A(t)}.
\]
Its expansion uses no additional progress induction.
\end{theorem}
\begin{proof}
Use external induction on the finite number of paired binders in $A$,
uniformly in all retained contexts. For $A(w)\equiv P(w)=Q(w)$, obtain
$P(s)=P(t)$ and $Q(s)=Q(t)$ by Theorem~\ref{zcc:numerical}, then use
root equality algebra to join them with $P(s)=Q(s)$.

For $A(w)=\mathsf D_{x,y}B(w,x,y)$, eliminate $A(s)$ at a fresh $a$.
The retained resource is $\mathsf Q_yB(s,a,y)$. Introduce the desired
$A(t)$ by case introduction with witness $a$ and fresh $z$. Its local
assumption is $B(t,a,0)$. Symmetry gives $t=s$, so the external induction
hypothesis transports that assumption backwards to $B(s,a,0)$.
QE then supplies $B(s,a,Sz)$. A forward transport by the same external
induction hypothesis gives $B(t,a,Sz)$. Close case introduction and DE.
Both recursive calls concern the strictly smaller matrix $B$. Freshness
conditions are met by renaming. This is structural reasoning about
finite formulas, not compound arithmetic induction.
\end{proof}

\begin{corollary}[Identity-language equivalence]\label{zcc:equivalence}
Replacing the supplied unrestricted counter-congruence rule by ZCC leaves
every native sequent derivable exactly when it was derivable before.
The same holds when comparing with the supplied unrestricted
object-transport calculus. Therefore the inherited whole-domain
bi-interpretation and equiconsistency with full first-order iPRA transfer.
\end{corollary}
\begin{proof}
Every ZCC inference is an old counter-congruence inference. Conversely,
Theorem~\ref{zcc:numerical}, with context $R(x,a;u,v.h)$, recovers every
old counter-congruence instance, including those with a state-dependent
body and arbitrary seed. Theorem~\ref{zcc:formula} recovers full object
transport. In the other direction, RL is full transport at the left
root, while ZCC follows by transporting one occurrence in a reflexivity
equation. The transformations recursively preserve end sequents and
freshness. The arithmetic and logical iPRA bridge is the attributed
supplied bridge; this refinement introduces no new domain or equality.
\end{proof}

\subsection{A different rule with only five nodes of overhead}

Define a separate primitive rule
\[
 \frac{\Gamma\vdash s=t}
 {\Gamma\vdash R(Ss,0;u,v.h(u,\vec p))=
                  R(St,0;u,v.h(u,\vec p))}
 \quad(\mathsf{ZCC}^{+}),
 \label{zcc:wrapped}
\]
again with no free occurrence of $v$ in $h$. Its premise endpoints are
arbitrary terms; the rule itself wraps them in successors. This is a
changed rule, \emph{not} ZCC with a side condition on its premise.

\begin{proposition}\label{zcc:wrapped-theorem}
Replacing ZCC by $\mathsf{ZCC}^{+}$ also gives the full native consequence
relation. Arbitrary numerical congruence now has a $P+5$-node construction
from a $P$-node source equality, without cases or progress induction.
\end{proposition}
\begin{proof}
For arbitrary $F$, apply $\mathsf{ZCC}^{+}$ using body $F(u,\vec p)$.
This directly gives $V_F(s)=V_F(t)$. Two RS endpoint computations and
two RL nodes give $F(s)=F(t)$, adding $1+1+1+2=5$ nodes.
Formula lifting is unchanged. Conversely, Theorem~\ref{zcc:numerical}
in $\mathcal C_0$ derives $\mathsf{ZCC}^{+}$ by choosing context
$R(Sx,0;u,v.h)$. Thus the two full calculi are equivalent.
\end{proof}

This factors numerical congruence through a fixed successor-computation
shape. It is not claimed to remove the cost of numerical binding or
substitution: that cost remains in the computation rules. No source-size
figure is assigned to this alternative calculus.

\subsection{Why successor-only premise endpoints are insufficient}

Consider instead retaining ZCC only when its \emph{premise} is syntactically
$Sp=Sq$, so its conclusion is $R(Sp,0;h)=R(Sq,0;h)$. Keep all other rules
unchanged and supply no rule wrapping arbitrary premise endpoints.

\begin{proposition}[A rule countermodel]\label{zcc:countermodel}
This successor-premise-only restriction does not derive full numerical
congruence and does not preserve the old native consequence relation.
In particular it does not derive
\[
                         0=S0\ \vdash\ S0=S(S0).
\]
\end{proposition}
\begin{proof}
Use the standard numerical domain $\mathbb N$. Interpret $0,S,R$ by
the ordinary zero, successor, and successor-seeded primitive recursion.
Interpret a native equation $s=t$ by
\[
 e(s^{\mathbb N},t^{\mathbb N}),
 \qquad e(m,n)\iff m=n\ \text{or}\ m,n\in\{0,1\}.
\]
For avoidance of ambiguity, this is a \emph{rule interpretation} of the
native equation atom; it is not asserted to be a model of ordinary
first-order logic with that atom as actual equality. Interpret resources
and paired formulas by their ordinary classical meanings
\[
 (\mathsf Q_yB)^\circ=B^\circ(0)\to\forall yB^\circ(y),\qquad
 (\mathsf D_{x,y}B)^\circ=
 \exists x\bigl(B^\circ(x,0)\to\forall yB^\circ(x,y)\bigr).
\]
Classical truth suffices for a rule-soundness countermodel of derivability.

The relation $e$ is an equivalence relation, so RL is sound. Every
computation equation has actually equal values and therefore is $e$-true.
For the restricted congruence rule, both counter values are positive.
On positive numbers $e$ agrees with actual equality. Consequently its
true premise implies identical actual counters, and the recursor
values agree. In fact this argument would allow arbitrary seeds and
updates whenever both premise counters are successors.

Hypothesis use, resource lookup, QE and DE are sound by their displayed
meanings. Case introduction is sound by standard numerical coverage.
Progress introduction is sound because ordinary induction on standard
$\mathbb N$ holds for the interpreted predicate, including predicates
whose equation atoms use $e$; only its actual truth set is relevant.
Freshness and substitution have their usual numerical meanings.
Thus every rule of the restricted calculus preserves truth in this
interpretation.

The premise $0=S0$ is true since $e(0,1)$ holds, while the conclusion is
false since $e(1,2)$ does not. Hence that sequent is underivable. The old
calculus derives it by successor congruence, proving the claimed strict
loss of native consequences.
\end{proof}

This countermodel rules out the proposed identity-language restriction.
It is not a theorem excluding every possible alternate interpretation,
changed equality convention, or enlarged axiom system.

\subsection{Evidence, overlap, and limits}

The source screen recursively inspected 32 distinct ZIP payloads from the
17 current archives, with 867 distinct text files and 443 parsed Python
ASTs. It found the expected old state-free-context remark but no explicit
zero-seed congruence restriction, wrapped rule, or successor-premise
countermodel. Pattern screening does not decide semantic duplication or
establish literature-wide priority. The result is attributed as a
refinement of the existing counter-congruence argument, not an invention
of its arithmetic or logical basis.

The supporting script \texttt{validate\_zero\_seed.py} preserves the
unchanged upstream readable checker and intercepts every CC call to
require literal-zero seed and absence of the outer state slot, correctly
tracking nested numerical binders. It constructs three symbolic proofs
(successor, initializer, and nested-update congruence), each with a one-node
source and exactly 21 nodes. A separate, explicitly coded wrapped-rule
branch verifies the corresponding three six-node proofs. All six pass.
Controls with a nonzero seed, a direct state occurrence, and a state
occurrence below another recursor are rejected. A finite check of the
countermodel relation is supplied only as a sanity check; the universal
rule-soundness argument is the preceding proof.

These artifacts verify concrete proof trees and side conditions. They are
not Lean/Rocq formalizations, universal correctness proofs of a new checker,
or measurements below the 558-byte current frontier. The numerical proof
bounds do not bound the size of formula-transport lifting, serialized term
labels, or runtime. No absolute simplest axiomatization or significant
checker-size breakthrough follows from this appendix alone.
